\chapter{為什麼現在所有寫作聽起來都一樣}

這場講座沒有真正的黑板數學。沒有黑板上的推導，沒有符號操作，也沒有投影片中可供還原的公式。然而，整場論證仍然有一道數學性的脊骨。Macfarlane 一再追問：某一種寫作形式，究竟如何回應某一種經驗形式；而主持人則不斷把這個問題磨得更尖，直到它變得局部、精確。於是，整場講座便透過區分、階梯、管線、尺度，以及能動性的轉換而展開。它從傳記性經驗開始，轉向光，進入散文的細微機械，繼而擴展到過程與世界觀，最後以修訂、精確，以及為獨特性辯護、抵抗扁平化而收束。

\begin{remark}
本章中的公式，都是對講座中明白說出的結構所做的示意性表述。它們不是從課堂記號中還原出來的。我們只在它們能澄清講座形式邏輯的地方使用它們。
\end{remark}

\section{群山、光，以及無法捕捉的失敗}

講座的開頭落在正確的位置：不是從教條開始，而是從讚嘆與起源開始。主持人問，Macfarlane 那種敏感性究竟從何而來；Macfarlane 則用訓練而非天賦的語言回答。他在群山之中長大。山會使感官變得銳利，要求警覺，也磨損掉平日自我的那層外殼。因此，危險並不是偶然的背景，而是知覺的發動機之一。

講座的第一條因果鏈，可以示意地寫成
\begin{equation}
\text{危險} \longrightarrow \text{警覺} \longrightarrow \text{開放} \longrightarrow \text{注意}.
\end{equation}
這裡已經不只是自傳了。這是一份關於作家如何被塑造的說明。更亮的光、更銳利的雪、像鐵絲一般刺進鼻中的空氣、風險評估、身體暴露：這些全都變成了注意的方式。山不只是供應風景，它供應的是一種感官教育。

從這個開場出發，主持人做出第一個決定性的轉折：``Tell me about that obsession with light.'' 他講起一個新的軼事：曾試圖從飛機上拍下落日，卻失敗了；隨之而來的是一種悲哀，因為他知道那個場景會在記憶中褪色，甚至也許永遠無法被完整說出。這種悲哀很重要。若跳過它，我們就會失去整場講座真正的驅動問題。Macfarlane 並不是從抽象美學出發，而是從一種切身感受到的無能為力出發：語言面對某一種看見的強度時，顯得不敷使用。

他的回答分兩步走。首先，他在言語中強化那道落日的力量：正如他所說，裡面幾乎有一種屠殺般的氣息，血色而狂野。接著，他退一步，說出那個形式性的教訓：當主題是光時，語言永遠來得太晚。

講座的核心對立就在這裡出現：
\begin{align}
\text{對應} &\not\Rightarrow \text{足以寫出光的文字},\\
\text{人工造作} &\Rightarrow \text{喚起知覺}.
\end{align}
如果我們死抓著「對應」不放，就會困在一場徒勞的追求裡：企圖讓語言與事物一點不差地貼合。Macfarlane 的建議更激進，也更解放：放棄這個追求。一旦我們不再堅持詞語必須逐點重現場景，隱喻、變形與形式構造便能開始做它們真正的工作。

\begin{definition}
講座所說的 \emph{對知覺的表述}，指的不是事物自身，而是事物如何經由一次特定的遭遇而抵達：它經過移動、心境、歷史與注意力的濾鏡。
\end{definition}

這也就是為什麼主持人那個印象派的類比如此貼切。Monet 並不是想把場景囚禁在顏料之中；他畫的是場景留在人身上的印象。Macfarlane 在散文裡做的是同一個動作。他甚至抗拒 ``capture'' 這個動詞，因為說「捕捉自然」，就已經把自然想像成俘虜了。更好的動詞是 evoke。一旦這個區分被安置好，講座便能從光，轉到另一種更明顯把形式與運動綁在一起的媒介：水。

\section{水、標點，以及思想的微小機械}

水並不是作為一個新題目進場的，而是作為反對「對應論」方法的第一個持續性測試。Macfarlane 花了多年寫一本關於河流的書，而從這場漫長勞作裡，他得出一個很強的形式結論：水沒有單一語法。整節的其餘部分，都繞著這一點旋轉。

主持人在這裡很有幫助，因為他立刻讓抽象落地。河流不是都以同一種方式流動。它們可以靜止、積蓄、蜿蜒、筆直、湍急。主持人一旦說出這些，Macfarlane 就能把這堂課磨得更銳利：
\begin{equation}
\text{河流的語言} \neq \text{只有快速的語言}.
\end{equation}
這個修正非常重要。若 ``流動'' 這個隱喻把一切水都壓扁成同一種散文質地，那它就太薄弱了。平展的湖水、轉折的河水，以及狂暴的急流，需要不同的節奏、語氣與句子結構。

按照講座自身的推進，河流分類隨即變成散文分類：
\begin{itemize}
\item 積蓄或靜止的水，要求延宕、開闊與停頓，
\item 蜿蜒的水，允許漂移、轉折與柔和的重定向，
\item 筆直的水，允許推進與線性，
\item 急流則要求迫切、累積與衝擊。
\end{itemize}

接著，講座便從大形式收窄到微小操作元件。Macfarlane 說自己既是介系詞狂熱者，也是標點狂熱者。首先出場的是破折號。句號，他說，會重重落下一個硬終點；破折號則保持液態。意義可以沿著它向後回旋，也可以沿著它向前流去。這不是一種裝飾，而是一種句法如何在不封死運動的前提下管理運動的模型。

接下來是介系詞，而在這裡，講座變得異常精確。介系詞不是填充物，它會改變作者與世界之間整個形上關係：
\begin{equation}
\text{關於河流} \longrightarrow \text{與河流一起} \longrightarrow \text{被河流書寫}.
\end{equation}
寫 \emph{about} 一條河，是把它當成客體來描述。寫 \emph{with} 一條河，已經是在思想中給了它某種共同作者身分，至少也給了它一種共同在場。而被寫於 \emph{by} river 之中，則更往前一步。作者變成導體、通道、媒介。主持人立刻就抓住了這個區分的力量，並用更平白的方式重述：幾乎就像河流成了共同作者。這種重述應該保留在筆記裡，因為講座就是反覆以這種方式向前推進的：Macfarlane 給出區分的壓力，而主持人則把它壓縮成一句更具教學性的句子。

下一步的收窄，則是節奏。詩教我們去期待節奏，但 Macfarlane 認為，非虛構作品其實也同樣作用於讀者所謂的 ``mind's ear''。它並不只靠命題運作，它也靠聲音、模式、節拍與內部押韻運作。這也就是為什麼他對第一句如此執著。第一句必須同時做幾件事，而講座實際上給了我們一個簡單模型：
\begin{equation}
\text{第一句} = \text{聲音模式} + \text{謎團} + \text{概念開口}.
\end{equation}

``The wind was rising, so I went to the wood'' 是講座裡最快樂的一個例子。它同時給了我們頭韻、節奏與謎團。風既然正在升起，為什麼反而要走進林中？這不只是漂亮措辭，而是一個小小的敘事引擎。相較之下，``12{,}000 years ago, a river is born'' 正如講座所說，節奏性較弱，卻更陰森不祥。它拿聲音上的鮮明，去交換時間尺度的巨大與概念上的神祕感。

然後，Macfarlane 又更細地推入內部聲音模式。``White as eyes'' 與 ``rises and flows'' 開始彼此回音；在任何論證還沒鋪開之前，文本就已經先在讀者身上發生作用。這一點是講座的重要提醒：非虛構文字並不會因為是非虛構，就只能平平地靠命題運作。

到了這裡，地基已經打好。水已經變成形式原理，標點與介系詞已被提升為真正的操作元件，而節奏也重新回到非虛構之中。於是，主持人就能提出整場講座最清楚的局部技藝問題。

\section{把速度寫出來，作為一個形式問題}

\subsection{Question \& Answer}

\paragraph{Question.} 我們怎樣寫出速度？

\paragraph{Answer.} 不是靠一條粗糙規則，例如「把句子寫短」。講座給出的答案更精確：我們是透過建立一種會翻滾、收縮，並且拒絕給讀者輕易停靠之處的句法，來寫出速度。

Macfarlane 做了主持人有資格要求他做的事：他拿出一個例子。不僅如此，講座還把這個例子本身放上螢幕，這一點本身就透露出一種教學方式。他不是在抽象地描述快速散文，而是停下來檢視一個標本。

\begin{figure}[ht]
\centering
\includegraphics[width=0.92\textwidth]{lecture_17_figure_02.png}
\caption{螢幕上的河流段落，展示一種速度與驚惶的句法}
\end{figure}

這張截圖本身就是證據。裡面沒有可見的方程式，也沒有講座圖解；那一整塊散文本身就是展品。在理論尚未被說明之前，眼睛便先感受到壓力：句群緊密、累積，幾乎不曾中斷。

Macfarlane 的解釋可以緊湊地寫成
\begin{equation}
\#(\text{and}) \uparrow,\qquad \ell_{\text{clause}} \downarrow,\qquad a_{\text{verb}} \uparrow
\Longrightarrow \text{驚惶與速度} \uparrow.
\end{equation}
其中 \(\ell_{\text{clause}}\) 是子句平均長度，而 \(a_{\text{verb}}\) 則代表愈來愈主動的動詞所施加的壓力。

他同時也提出一個至關重要的局部等價關係：
\begin{equation}
\text{``and''} \sim \text{破折號}.
\end{equation}
也就是說，在這一類散文裡，連接詞承接了先前破折號所做的工作。它讓語句繼續流下去。它不讓平靜的層級出現，也拒絕那種由句號與整齊從屬關係所構成的停起從容。

這個推演很簡單，卻很準確。我們可以照著講座本身的順序把它寫出來：
\begin{enumerate}
\item 從一個已經在移動中的句子裡開始，而不是從邊界清楚的子句起步。
\item 重複使用 ``and''，讓一個單位不斷向下一個單位溢流。
\item 讓連接詞之間的子句變得更短、更快。
\item 提高動詞的壓力：slammed、exploded、punching、ramming、hauling。
\item 在 ``river'' 前面去掉冠詞，於是這股力量不再像一個被命名的物件，而更像一種無所不在地作用的媒介。
\item 結果便是：讀者開始喘不過氣，因為句法拿走了那些通常會恢復呼吸與層級的停點。
\end{enumerate}

去掉冠詞，尤其是一個很有力的小動作。``The river'' 是一樣我們可以與之保持距離的東西；``river'' 則更像一種能動性、一種撞擊、一種不是風景而是媒質的存在。看似微不足道的語法扣除，改變了整個場景的本體論。

由於螢幕上只呈現那一整塊散文，而沒有任何分析圖，我們若要補圖，就必須清楚讓它維持在次要位置。因此，下圖只是編者的重建，而不是還原出來的講座圖像。

\begin{figure}[ht]
\centering
\begin{tikzpicture}[x=1cm,y=1cm]
\draw[rounded corners] (-5.2,1.2) rectangle (-3.6,2.0);
\node at (-4.4,1.6) {and};
\draw[rounded corners] (-2.9,1.2) rectangle (-1.3,2.0);
\node at (-2.1,1.6) {and};
\draw[rounded corners] (-0.6,1.2) rectangle (1.0,2.0);
\node at (0.2,1.6) {and};
\draw[rounded corners] (1.7,1.2) rectangle (3.3,2.0);
\node at (2.5,1.6) {and};

\draw[->] (-3.6,1.6) -- (-2.9,1.6);
\draw[->] (-1.3,1.6) -- (-0.6,1.6);
\draw[->] (1.0,1.6) -- (1.7,1.6);

\draw[rounded corners] (-5.7,-0.2) rectangle (-3.1,0.7);
\node[align=center] at (-4.4,0.25) {子句\\縮短};

\draw[rounded corners] (-2.8,-0.2) rectangle (-0.2,0.7);
\node[align=center] at (-1.5,0.25) {動詞變得\\更主動};

\draw[rounded corners] (0.1,-0.2) rectangle (3.0,0.7);
\node[align=center] at (1.55,0.25) {``river''\\而非 ``the river''};

\draw[rounded corners] (-2.0,-1.8) rectangle (1.0,-0.9);
\node[align=center] at (-0.5,-1.35) {讀者的\\喘不過氣};

\draw[->] (-4.4,1.2) -- (-4.4,0.7);
\draw[->] (-2.1,1.2) -- (-1.8,0.7);
\draw[->] (0.2,1.2) -- (1.2,0.7);

\draw[->] (-4.4,-0.2) -- (-1.05,-0.9);
\draw[->] (-1.5,-0.2) -- (-0.7,-0.9);
\draw[->] (1.55,-0.2) -- (0.0,-0.9);
\end{tikzpicture}
\caption{對螢幕上河流段落背後機械的編者式重建}
\end{figure}

講座對這整體效果所用的說法極好，應當保留：這是一種 ``syntax of panic''。成就不在於句子 \emph{描寫} 了驚惶，而在於它讓讀者經歷一種受控的、與驚惶相似的感受。

到這裡，主持人做出了下一步最正確的推進。若一個句子可以被這樣建造，那麼這類句子究竟是從哪裡來的？

\section{從筆記本到一本書}

\subsection{Question \& Answer}

\paragraph{Question.} 筆記如何變成一本書？

\paragraph{Answer.} 不是因為一本完成的小書早就藏在筆記本裡，只待被擴寫出來。講座堅持，筆記一開始是碎片化的、加密的，往往甚至不可交流。一本書並不是那些筆記本來就已經是的東西；一本書是它們在反覆回返、重新喚回與模式辨識之下，漸漸成為的東西。

講座的節奏在這裡改變了。經過河流句子的迫切性之後，Macfarlane 故意放慢，幾乎是在說：讓我把它拆成幾個階段。這個步調很重要。主持人已經把問題從句子機械推到書本機械，而 Macfarlane 的回答，也因此從警句轉向過程。

第一個階段是田野工作。這些書要花上數年，需要多次旅行，也需要與人、河流、山脈反覆相遇。Macfarlane 偏好使用小筆記本，大約是 A6 尺寸，因為它們可以隨身攜帶、彎折，也能在真實條件下被使用。被記進去的，是他半開玩笑所稱的 \emph{qualia}：零星碎片、片語、對話殘片、圖像閃光，以及那些黏在心上的小石子與羽毛。

這些筆記並不整齊，而是斷續的。這是整場講座最有力的實際觀點之一。當遭遇正在發生時，根本沒有餘裕去寫成熟的散文。人只能先把東西記下來，然後繼續往前。每天收工時的那種補記也在這裡變得重要：額燈開著，把仍然熱乎乎的記憶從腦中拉下來，趁它還沒冷掉。

接著，主持人提出一個極好的隱喻：頁面成了一個凝結面。思想原本像蒸氣一樣在場，而紙張則讓它開始凝聚。Macfarlane 接住這個比喻，並把它說得更強：紙與筆確實會做出鍵盤與像素做不到的事。那些筆記是 ``fresh from the fire'' 的，仍然帶著熱度，還沒有被塑形成另一位讀者可以消化的樣子。

我們可以把「筆記本到書」的管線寫成
\begin{equation}
\text{碎片} \longrightarrow \text{記憶觸發器} \longrightarrow \text{記憶擴展} \longrightarrow \text{模式辨識} \longrightarrow \text{心之工作}.
\end{equation}

這條鏈的中段尤其值得強調。當 Macfarlane 回到家，在螢幕上重讀筆記本時，每一個碎片都成了某條線頭可見的一端。只要拉動那條線，整個周遭場景便又重新打開。這是講座對回憶現象最精確的描述。筆記不是經驗本身，而是讓人重新進入經驗的把手。

到了這裡，Macfarlane 提到了 Rilke，雖然非常謹慎，而且只是轉述。逐字稿在這一段不穩定，但區分本身很清楚：
\begin{equation}
\text{筆記本工作} \longrightarrow \text{心之工作}.
\end{equation}
即時的記錄並不是最終的勞作。真正困難的工作開始於：人回頭去看那些圖像，把它們彼此關聯起來，並理解它們之間的共鳴。

這也就是為什麼講座接著轉向那隻張開的手。這個例子之所以出色，在於它不是事後任意挑選出來的。Macfarlane 最初是在洞穴模板手印藝術中注意到那隻手：不是手本身的痕跡，而是它缺席所留下的痕。這個圖像後來又在其他旅程、其他媒介裡反覆出現，甚至還在社會性的形式裡再現，例如問候與歡迎。一個起初孤立的圖像，於是變成了一個母題。

我們可以更精確地把這種浮現寫成
\begin{equation}
\text{反覆出現的圖像} \longrightarrow \text{跨旅程的重現} \longrightarrow \text{模式點亮}.
\end{equation}
這裡的關鍵詞是 ``lighting up''。母題並不是單純重複，而是開始活化。作者看見了它；讀者也因而被賦予了看見它的能力。

主持人在這裡很有用地把這件事重新說了一遍：所以你真正描述的，其實就是模式辨識。正是如此。而正是這種辨識，才把一堆加密碎筆，轉換成一本有形可辨之書的開端。

講座接著又補上它最不浪漫、卻也最應保留的一條建議，而且這建議正是從模式討論自然長出來的：屁股坐上椅子。出現。第一條規則是做工；第二條規則是尋找模式。兩者不能分開。

\section{驚奇、尺度與去學習}

\subsection{Question \& Answer}

\paragraph{Question.} 去學習，究竟是什麼意思？

\paragraph{Answer.} 講座只有在建立了句子層面的技藝與長篇過程之後，才走到這個問題。這個順序很重要。去學習不是作為某種神祕裝飾被提出的，而是作為一位已經學會注意、學會辨識模式，也已經學會懷疑天真對應論的作者，所將抵達的更大認識論結果。

眼前直接的開場，是 Macfarlane 講到自己的兒子。他說自己正在寫一本名為 \emph{Is a River Alive?} 的書，而孩子幾乎是立刻回答：``Of course.'' 主持人立刻感受到其中的力量。成年人常常失去一種既清晰、又帶著魔法感的世界關係；孩子也許不一定有更好的論證，但卻保有一種較少被收窄的本體論。

Macfarlane 的第一個答案，是驚奇。驚奇，他說，是一種必要的生存技能。彩虹於是成了講座的標本。這個要點同時朝兩個方向成立。第一，驚奇是面對這個自由贈與的世界時，張口結舌的訝異。第二，解釋並不會取消驚奇。彩虹可以被描述成：光在水中被分解為不同波長；但這個解釋並不會溶解那份經驗。相反地，講座暗示，科學有時甚至能更細膩地把現實推向驚奇，而不是把驚奇剝除。

接著，Blake 出場了；隨著 Blake 一起進場的，是尺度。一粒沙可以打開一個世界，一首極小的詩也可以包含永恆的預感。從這裡，Macfarlane 繼續概括：
\begin{equation}
\text{微觀尺度} < \text{人類尺度} < \text{宏觀尺度}.
\end{equation}
這當然只是示意性記號。講座真正的意思，不是嚴格排序，而是並置與嵌套。我們平常住在人類尺度之中，但作者可以把這個尺度放到顯微與宇宙的尺度旁邊，並讓它們之間的張力保持可見。

這也就是為什麼講座會提到近代早期透過顯微鏡與望遠鏡而被撬開的尺度世界。Galileo 看見月球上的山脈；van Leeuwenhoek 看見水滴中的生命。對 Macfarlane 而言，寫作可以完成一種類似的動作：把不同尺度並置，而不把它們壓扁成一個。

直到這時，主持人才引入 ``blinders of rationality'' 這個說法。這正是恰當的轉場。Macfarlane 立刻同意：對那些在理性主義習慣中受訓的人來說，要相信河流以某種超過其中各個生物總和的方式活著，確實需要去學習。他接著給出講座裡最好的圖像之一：理性主義不是一盞照亮整個洞穴的探照燈，而是黑暗中一塊移動的光斑。我們所知道的，總是比那塊光斑誘惑我們以為自己知道的要少。

那段河上旅程，則為這個哲學轉向提供了物理基礎。講座在這裡短暫變得帶有數值色彩：
\begin{align}
Q_{\text{normal}} &\approx 150\,\mathrm{m^3\,s^{-1}},\\
Q_{\text{observed}} &\approx 275\,\mathrm{m^3\,s^{-1}},\\
\frac{Q_{\text{observed}}}{Q_{\text{normal}}} &\approx 1.83.
\end{align}
也就是說，那條河的流量幾乎達到了平常的兩倍。這裡的算術很簡單，但講座用得很好。它把形上主張錨定在肉身磨難之中。這並不是在象徵性的水邊進行安靜沉思，而是長時間的翻滾、游泳、撞擊、精疲力竭，以及日復一日的磨蝕。

Macfarlane 的說法令人難忘：是河流把他去學習了。這種去學習既是身體性的，也是形上學性的。它打破了那個以為「這不過是水」、不過是死物質的假設。河流開始以能動性、在場、力量與意志的形式呈現自己。這種經驗在旅程終點達到高潮：一種河流存在，甚至近乎神性的在場。講座在這裡很小心。語言無法完整承載整個事件。它最多只能註冊某種與所發生之事相似的痕跡。

而這，又把我們帶回先前的方法。即使在這裡，也許尤其是在這裡，目標仍不是完美表述，而是在壓力之下，精確登錄知覺。

\section{問題作為入口，而他人作為共同作者}

\subsection{Question \& Answer}

\paragraph{Question.} 一個簡單的問題，怎麼可能打開整整一本書？

\paragraph{Answer.} 因為真正好的問題，並不是一條等待答案的口號；它是一個狹窄的開口，通向一個比字面上大得多的房間。

主持人此時點出了 Macfarlane 這些書的一個結構性特徵。它們常常從一些看似狹窄的問題出發：Why climb mountains? Can a forest think? Does a mountain remember? Is a river alive? Macfarlane 同意，但在解釋結構之前，他先停下來，談一個失敗的答案：George Mallory 的 ``Because it's there.'' 它當然有名，但卻並不好用。對一本書而言，它在應該保持開放的地方，把探問提前封死了。

於是，講座的「入口模型」便隨之而來：
\begin{equation}
\text{簡單問題} \longrightarrow \text{狹窄入口} \longrightarrow \text{龐大的複雜之室}.
\end{equation}
洞穴探勘的例子，給了我們最恰當的物理直覺。人得先擠過一個小入口，也許還得潛過一道積水低點，然後才進入一個巨大的內部空間。主持人接著又用 Narnia 的衣櫥類比，把同一點以另一種語域說得更清楚：衣櫥很小，後面那個世界卻不小。

這不是事後附會的比喻。Macfarlane 說，這實際上幫助他把這些書的問題究竟是什麼樣的東西，結晶出來。它們就是入口：門檻上看來不大，越過去之後卻極為寬闊。

接著，講座又回到實際層面。一旦那個入口問題被找到，Macfarlane 便會寫一封信給未來的自己。因為這些書要花很多年，寫書的人與世界本身，都會在書成形的過程中改變。因此，這封信不是誡命，而是一份暫時性的希望聲明：這本書也許會變成什麼，它可能會反覆出現哪些隱喻，它可能會在世界中做什麼工作。

還有另一種輔助：釘在牆上的引文。談到 \emph{Is a River Alive?}，講座回到 Ursula Le Guin，尤其回到那句 ``a great reach outward of mind and imagination.'' 這句話成了工作中的指令。問題打開門，而這種向外伸展，則讓人有能力穿過那扇門。

接著，主持人問：寫一本書究竟有多少是計畫，又有多少是冒險？Macfarlane 用另一個河流類比作答，而這個類比太好了，不應被壓扁。這本書並不像乘著筏子順流而下，而更像是在河裡逆流步行。雙腳都踩穩時，人是穩定的。真正的前進，恰恰開始於你抬起一隻腳的那一刻，因為也就在那一刻，水流有能力把你沖失平衡。這個教訓同時是實際的，也是概念性的。一本書不會只因為存在而使人失穩；它是在逼你跨出下一步時，才真正讓你失去平衡。

這個回答自然引向他人。Macfarlane 抗拒那種較舊、排除人的自然書寫。他希望書裡充滿同行者、專家、在地知者，以及令人驚嘆的聲音。主持人立刻看出其中的力量，於是問：這些人幾乎像是共同作者嗎？Macfarlane 以精神上的意義回答說，是的；而講座中的例子也完全證明了這一點。恰當的人，會把孤身作者永遠無法提供的知覺形式帶進書裡。

此時，主持人對對話的隱喻變得關鍵：一起生火。這是對的，因為對話不只是交換內容，它也會產生兩個說話者各自都無法單獨製造的熱與光。Macfarlane 很高興把這個模型直接引進寫作。

講座於是再次擴展，這回是橫跨形式。詩、歌劇腳本、歌曲寫作與合唱寫作，都在教長篇散文作者一些單靠散文本身可能會遮蔽的東西：可歌唱性、聲流、詞與詞之間的阻力、圖像得以稍微鬆散共存的方式，以及輕微錯位的形式用途。歌曲寫作更讓他想到了 glitch 的概念：不是錯誤，而是一種有生產力的摩擦。踩在拍子外的設置、令人不安的滑移、以及尚未解決的接觸，都可能成為形式資源，尤其當作者想讓語言註冊陌異，而不是把陌異磨平時。

這便為講座最後一個後段動作做好準備。當一本書已經有了入口問題、寫給未來自己的信、合作者，以及更寬的形式資源，技藝上的問題便變得更直接：我們怎樣在不被推回平均散文的前提下，保住精確、修訂、身體力道與個體性？

\section{精確、修訂與獨特性}

主持人在這一段的轉入方式恰到好處：``We've got to talk about words.'' 講座此時進入一段快速的後程奔跑，但它並不零散。它其實是一個聚攏起來的、對書名問題的實際回答。

Macfarlane 先從詞彙表與地方詞說起。他使用不尋常的詞，不是為了聽起來不尋常，而是因為它們精確。他堅持，精確不是抒情性的敵人，而是抒情性的工具之一。一個單一的地方詞，可能就壓縮了整整一段知覺。古英語的 kenning，像 ``bone-cage'' 與 ``whale-road''，展現了語言如何把隱喻與命名熔於一爐。Seamus Heaney 那句 ``the palp and heft of language''，指向的也是同一種觸覺智慧：詞語可以被捏、被掂量、在口中與手裡被感到。

接著，講座轉向 Gaelic 地名與地方詞。這一段逐字稿中的若干拼寫顯然不穩，精確正寫應在後續另行核對。不過，其中的基本意思是穩固的。許多地名本身編碼了形狀、方位與景觀關係。它們不是貼在先存地形上的標籤，而是注意力的儲藏所。講座半開玩笑地把這稱作一種 ``Gaelic positioning system''，但玩笑裡其實帶著真正的形式洞見：語言可以儲存地形學。

於是，語言死亡也進場了，而講座對它下的定義異常清楚。一種語言會死亡，是當它最後一位活著的說話者去世，而那種語言不再被傳遞下去之時。記錄也許仍在，但傳遞停止了。隨後更強的主張立刻跟上：
\begin{equation}
\text{語言死亡} \longrightarrow \text{知識流失}.
\end{equation}
因為語言本身就是一套知識儲存系統，所以有些知識會隨之消失，尤其是在那些知識無法被乾淨地譯入另一種語言的地方。於是，語言死亡同時是一種文化損失，也是一種生態損失。

講座接著把問題拉回英語本身。Macfarlane 的回答很細。某些時候，英語應當像玻璃一樣透明，讓讀者忘記媒介；另一些時候，它應該濃厚、豐饒、可觸，讓讀者感到每一個詞。若這兩種模式其中之一被誤判，散文便會成為干擾；若選擇得當，讀者的身體便會參與到語言的重量或抬升之中。這個想法把 Heaney 的 ``palp and heft'' 延伸成了一條一般性的技藝原理。不是每個句子都該宣告自己，也不是每個句子都該消失不見。

接著進入修訂，而在這裡，Macfarlane 給出了講座裡兩個最有力的實作模型。第一個是陶工模型：
\begin{align}
\text{濕黏土團} &\longrightarrow \text{旋轉} \longrightarrow \text{塑形} \longrightarrow \text{裝飾},\\
\text{原始草稿} &\longrightarrow \text{反覆多次} \longrightarrow \text{成形} \longrightarrow \text{完成}.
\end{align}
這不只是裝飾性的比喻，而是一種對「完成的散文會乾淨地一次到來」這種幻想的拒絕。最初的材料是泥濘的，形式則是經由反覆塑造才浮現出來。

第二個模型則是馬賽克式的組裝：
\begin{equation}
\text{段落塊}_1 + \text{段落塊}_2 + \cdots + \text{段落塊}_n
\longrightarrow \text{透過組裝而成章}.
\end{equation}
Macfarlane 並不要求線性起草。若某一段被卡住，另一段完全可以先寫。一本書是從若干塊面拼起來，再被安排成流。與此同時，他還給了兩條很實際的建議：一天工作結束時，讓下一句已經被知道；若真的卡死，就往下游跳，先寫另一段。這兩者都是維持運動、而不是膜拜阻滯的方法。

講座反覆回來的一句口頭禪，在這裡又一次有力地出現：屁股坐上椅子。出現。即使一天只產出一段，那本書也還是多了一段。這也許是整場講座裡最不華麗的一句話，卻同時也是最有用的一句。

接下來，自動化登場了。Macfarlane 區分的是：這類工具對平均散文做了什麼，以及它們對有個性的散文做了什麼。文法檢查器確實改善了許多普通句子，同時也會把刻意的偏離標示成錯誤。他舉出的例子——無動詞段落、碎句、突兀分段、句法繃緊——恰恰就是那些自動平滑化最討厭的寫法。這裡的形式交換關係很容易寫清楚：
\begin{align}
\text{工具平滑化} \uparrow &\Longrightarrow \text{平均稱職度} \uparrow,\\
\text{工具平滑化} \uparrow &\Longrightarrow \text{獨特性} \downarrow.
\end{align}
這是講座在後段最清楚的一條主張之一。當寫作被往平均可接受度方向修補時，它就開始聽起來彼此一樣。

接著，主持人又重新打開另一個相關但不同的問題：我們怎樣讓語言具有 visceral 的效果？這個回答不是推翻，而是延續了前面對 ``capture'' 的拒絕。有身體感的寫作，並不是捕捉事件本身，而是把事件的一種身體類比轉移到讀者身上。在談到 \emph{Underland} 中那些令人幽閉的段落時，Macfarlane 借用了 William Golding 的說法：``sympathetic kinesthesia.'' 讀者的身體會收緊、加速、退縮。這並不是反形式的，而是形式力量在命題之下運作。

直到這一切都說完之後，主持人才回到那條關於刪去贅詞的名言。Macfarlane 的回答既有紀律，又不偏激。Walter Pater 的 ``hard, gem-like flame'' 指向一條有效的傳統；Hemingway 與 Carver 也站在類似的減法系譜上。但 Macfarlane 拒絕把減法普遍化。有些風景要求刪除，另一些則要求增加。這裡真正的關鍵詞是 bespoke。每一次遭遇，都要求它自己的密度、自己的開口、自己的節奏。

而這種判斷，在講座最後談到驚異時就變得更重要了。他說，誇飾是敵人。華麗濫情的散文會說謊。驚異最好的呈現方式，不是把語言不斷膨脹，直到它開始吼叫。更好的方式是精確、節制，以及拒絕過度解釋。主持人在這裡提供了一句極強的轉述：當我們把每一個 \emph{i} 都點上、每一個 \emph{t} 都劃上時，我們就不再給讀者留下任何可供凝視與在句中起身活動的空間。

而這，又把講座優雅地帶回入口的觀念。來自 \emph{Is a River Alive?} 的一句話問道：如果英語裡沒有一個動詞叫做 ``to river''，那又有什麼比一條河更像動詞？這句話既完成，又未完成。它讀來像是結束了，但同時又把人往前送。這正是講座希望好句子具備的行為方式：它既要給出形狀，也要留下神祕；它既要在局部上滿足，也要在整體上重新打開。

\section{Summary}

講座始於山，終於獨特性，而其中的推進卻極為緊密。危險的地景訓練出警覺；無法 ``capture'' 光，迫使人從對應論轉向人工造作；水則教會我們，散文必須回應的不只是一種運動，而是多種運動。標點與介系詞因此升格為真正的操作元件。速度則是透過串聯式並列、收縮的子句、主動的動詞，以及從 ``river'' 中拿掉冠詞來寫出的。筆記之所以能變成書，只能經過回返、記憶線頭、模式與心之工作。驚奇拓寬成一個關於尺度的寫作模型，然後又深化為更艱難的要求：去學習。一本書透過入口問題、寫給未來自己的信、合作者與更寬的形式資源而被組織起來。到了後段，寫作的技藝便成了一場防衛：為了保住精確、修訂、身體力道與個體性，去抵抗平滑化與同質化。

我們可以把講座最後的主張壓縮成一條示意關係：
\begin{equation}
\text{注意} + \text{形式} + \text{判斷} \Longrightarrow \text{獨特的散文}.
\end{equation}
這就是對書名的實際回答。當寫作失去它的耳朵、它的壓力，以及它與活生生知覺之間精確的關係時，它就會聽起來彼此一樣；而當句法、節奏與結構重新對那場召喚它們誕生的遭遇負責時，寫作便重新活了過來。